\documentclass[../Main.tex]{subfiles}

\begin{document}

\section{Reflexion} \label{sec:Reflexion}

\subsection{Reflexion des Arbeitsprozesses}

Dieses Projekt hat etwa ein Jahr unserer Zeit in Anspruch genommen und mehrere Phasen durchlaufen, die wir im Folgenden der Reihe nach aufgreifen. Diese Reflexion deckt sich mit den Einträgen in unseren Arbeitsjournalen, die wir fortlaufend verfasst haben.

Ganz am Anfang waren wir noch völlig unsicher und hatten ganz andere Vorstellungen davon, welches Produkt am Ende entstehen würde. Also stocherten wir in alle Bereiche hinein: Wir begannen mit Spielsystemen, von denen wir ohne Recherche annahmen, dass wir sie brauchen, überlegten, wie wir die Erkenntnisfrage beantworten wollen, und recherchierten, wie man recherchiert. Bis zur Zwischenpräsentation blieben uns noch vier Monate. Wir waren wenig zielorientiert und machten einfach, was uns gefiel oder was wir noch nicht ausprobiert hatten. Das mag planlos wirken, rückblickend war es aber ein Auskundschaften dessen, was es gibt. Die Fallhöhe war zudem gering: Was auch immer wir machten, nichts musste verworfen werden, und die Arbeitsschritte hingen noch kaum voneinander ab. Dass das aufging, war zwar eher ein glücklicher Zufall als strukturiertes Denken unsererseits, doch so ist die Natur eines Projektanfangs. Unter Leuten, die Projekte machen, heisst es ja, sie hätten nie angefangen, hätten sie gewusst, wie viel Arbeit es wird. Könnten wir noch einmal beginnen, würden wir es genauso angehen, mit dem einzigen Unterschied, dass wir wüssten, dass wir am Anfang eine Narrenfrist haben.

Bis zur Zwischenpräsentation spielte die Erkenntnisfrage eine untergeordnete Rolle, denn die dafür relevanten Systeme beruhten auf einer technischen Basis, die wir noch nicht hatten: einem Spielercharakter mit Animationen und Steuerung, Gegnern und Code, der das Spiel überhaupt zum Starten bringt. Obwohl wir dabei die Recherche vernachlässigten, war das rückblickend die richtige Entscheidung. Zum einen machte uns dieser Teil am meisten Freude, zum anderen mussten wir uns zuerst eine Vorstellung von unserem eigenen Spiel bilden. Erst nachdem wir die ersten Systeme gebaut hatten, wussten wir, worauf wir in anderen Spielen achten mussten, und konnten erst recht eine gute Recherche machen. Bei einem nächsten Projekt würden wir die Reihenfolge davon abhängig machen, wie konkret unsere Vorstellungen sind und wie gross unser Kenntnisstand ist: Haben wir beides, recherchieren wir zuerst. Fehlt eines davon, sammeln wir zuerst eigene Erfahrungen -- also genau das, was wir diesmal eher unbeabsichtigt getan haben.

Gleichzeitig rückte die Zwischenpräsentation näher. Features, die wir zwar wollten, die für die Erkenntnisfrage aber wenig relevant waren, etwa ein Stylemeter wie in \game{Devil May Cry}\autocite{capcom_devilmaycry_2001} oder eine zweite, auswechselbare Waffe, mussten wir schon jetzt streichen oder stark eingrenzen, und wir mussten priorisieren, was für ein Spiel essenziell ist. Dazu kam ein begründet geplanter Playtest, und die Checkliste aus dem Vertrag hatte noch viele offene Punkte. Wir haben alles erreicht, was wir vertraglich festgehalten hatten, auch wenn es sehr, sehr knapp wurde. Nicht das Fertigstellen, sondern die Qualitätssicherung hat uns einen Strich durch die Rechnung gemacht, weil wir möglichst wenige Kompromisse eingehen wollten. Im Nachhinein wäre es interessant gewesen zu erfahren, wie 80 oder 85 statt 100 Prozent unseres Anspruchs bewertet worden wären. Da die Zwischenbewertung nicht in die Note der Maturitätsarbeit einfliesst, ist das eine verpasste Chance. Andererseits konnten wir Erfahrung mit Crunchen sammeln, die uns am Ende ermöglicht hat, mit genügend Zeit und kaum Crunchen 95 bis 100 Prozent herauszuholen.

Zur Zwischenpräsentation erhielten wir eine Rückmeldung: Vieles war gut, einiges weniger. Die Kritik betraf weniger den Prototyp als alles rundherum, besonders die begleitenden Texte: Wir mussten unsere Aussagen so formulieren, dass erkennbar ist, ob es sich um eigene Einschätzungen oder um belegte Aussagen handelt, und unsere KI-Nutzung vollständig dokumentieren. Beim Umsetzen hatten wir zunächst Schwierigkeiten, denn hätten wir die bemängelten Punkte verinnerlicht gehabt, hätten wir sie von Anfang an richtig gemacht. Deshalb stellten wir unserer Betreuungsperson so lange Fragen, bis wir sicher waren, sie verstanden zu haben.

Nach der Zwischenpräsentation hatten wir wieder genug Zeit, und die Abgabe war weit genug entfernt, dass wir unsere Ambitionen beim Projektumfang vergrösserten. Gleichzeitig nahm das anfängliche Tempo ab, mit dem wir Neues produzierten. Zum einen war das Projekt inzwischen so gross, dass kleinere Neuerungen wenig Unterschied machten. Zum anderen war unser anfänglicher Eifer zwar weit davon entfernt, erloschen zu sein, doch das Projekt zu pflegen und Aufgaben zu erledigen, die uns weniger Spass machten und sich mittlerweile angehäuft hatten, setzte uns langsam zu. Wir nahmen deshalb beide etwas Abstand vom Projekt und verfolgten andere Interessen, was möglich war, weil wir gut vorgearbeitet hatten. Danach kehrten wir gestärkt zurück, zumal wir die nächste Deadline im Blick hatten und uns ohnehin keine andere Wahl blieb.

Mittlerweile hingen die einzelnen Arbeitsschritte so stark voneinander ab, dass wir in unserer Partnerarbeit nicht mehr einfach tun und lassen konnten, was wir wollten, sondern uns aufeinander abstimmen mussten. Wir arbeiteten deshalb an möglichst grossen, voneinander unabhängigen Blöcken und vereinbarten ein Datum, wenn sich unsere Arbeiten kreuzten. In dieser Zeit entstand auch fast die gesamte restliche Recherche. Das funktionierte sehr gut, wenn beide rechtzeitig fertig wurden, was allerdings nicht immer der Fall war. Künftig würden wir deshalb einen dritten, grösseren Arbeitsstrang offenhalten, an dem weiterarbeiten kann, wer gerade durch die andere Person blockiert ist.

Nach Abschluss der Recherche wussten wir genau, welche Systeme noch implementiert werden mussten, und mit der näher rückenden Abgabe hatten wir nun auch die Beantwortung der Erkenntnisfrage im Blick. Hier beginnt der wirklich gut geplante Teil unserer Arbeit. Da wir nun ein Fundament hatten, das in eine klare Richtung ging, und wussten, dass unsere Arbeit nicht mehr verschwendet sein würde, konnten wir zielgerichtet auf die Erkenntnisfrage hinarbeiten. Dieser Teil ist uns besonders gelungen, weil wir auf Hinweis unserer Betreuungsperson einen Plan bis zur Abgabe erstellt haben. Dass dieser Plan gut wurde, verdanken wir den Erfahrungen aus der Zwischenpräsentation, bei der wir bereits gelernt hatten, wie man auf eine Abgabe hinarbeitet und plant.

Unseren ersten fertigen Export hatten wir am 25.~Oktober. Technisch war das Spiel damit fertig, doch es spielte sich grausam, und einige Rechercheergebnisse hatten wir noch nicht konsequent umgesetzt. Es war allerdings halb so schlimm, wie es anfangs den Anschein hatte: Wir mussten nur an ein paar Stellschrauben nachjustieren und die fehlenden Aspekte ergänzen, und das Spiel wurde unserer Recherche gerecht und war bereit für die Playtests. Diese Erfahrung hat uns gezeigt, dass der Teufel im Detail steckt. Das Problem entstand genau dort, wo Recherche und Prototyp vollständig ineinandergreifen mussten -- sie griffen schlicht nicht gut ineinander. Bis dahin hatten wir Recherche und Prototyp getrennt betrachtet. Anfangs war dieser Gedanke nicht ganz falsch, doch je weiter wir arbeiteten, desto mehr hingen die beiden zusammen. Beim nächsten Mal wissen wir, dass solche Teile schon im Voraus gut aufeinander abgestimmt werden müssen und nicht erst als nachträglicher Gedanke.

Gegen Ende der Arbeit führten wir die Playtests durch. Sie verliefen recht gut, da wir bei der Analyse von \game{NieR: Automata} und \game{Ninja Gaiden} bereits selbst Spieletests durchgeführt hatten. In der Schlussphase bereiteten uns eher Formalitäten Sorgen. Auch diese Phase verlief unserer Einschätzung nach ziemlich gut, nicht weil wir darin etwas Besonderes gemacht hätten, sondern weil wir zuvor gut geplant und uns genug Zeit dafür gelassen hatten. Zudem war uns bewusst, dass es knapp wird, wenn wir uns nicht an den Plan halten, auch wenn es noch nach genug Zeit aussah.

\subsection{Gruppenarbeit und Aufgabenteilung}

\begin{itemize}[leftmargin=1.5cm]
    \item \textcolor{red}{\textbf{1. Wie ihr die Arbeit aufgeteilt habt.} Laut Zwischenbewertung: Rodin stärker bei Animation, UI, Bericht und README; Zhile stärker bei Architektur, Bewegung, State Machines und Kampf-Grundlagen. Das deckt sich mit Git: Zhile vor allem \code{Entities/Player}, \code{Common} (Utility, Tools), \code{Core}, Kamera; Rodin vor allem \code{Interfaces} und Menüs sowie die TrenchBroom-Karten des Levels.}
    \item \textcolor{red}{\textbf{2. Tabelle Aufgabenteilung und Eigenleistung} für Prototyp und Bericht (Rodin / Zhile / gemeinsam). Zum Bericht: Kapitel 1--8 und die \game{NieR}-Analyse Rodin, die \game{Ninja-Gaiden}-Analyse Zhile, README und Architekturdiagramm (Juni) vor allem Zhile. Bitte prüfen und ergänzen (Playtests, Leveldesign, Audio, \ldots).}
    \item \textcolor{red}{\textbf{3. Zusammenarbeit und Kommunikation:} Wie habt ihr euch abgesprochen (Treffen, Discord, GitHub)? Wie habt ihr Entscheidungen getroffen?}
    \item \textcolor{red}{\textbf{4. Veränderung im Verlauf:} z.\,B. ab Oktober Rodin vor allem Level, Playtests und Bericht, Zhile weiter am Spiel (laut Git im Oktober sehr aktiv). Wie kam es dazu?}
    \item \textcolor{red}{\textbf{5. Fazit:} Was lief gut, was war schwierig, was würdet ihr in einer Gruppenarbeit anders machen? Hinweis: Commit-Zahlen sagen wenig über den Aufwand aus, deshalb nicht als Massstab verwenden.}
\end{itemize}

\subsection{Verwendung von KI} \label{sec:ReflexionKI}

\game{Project-Polearm} war ein grosses Informatikprojekt, da muss doch viel KI im Spiel gewesen sein, oder? Eigentlich nicht. Unsere Philosophie war, dass KI nichts erschaffen soll, sondern nur ein Werkzeug ist, das besonders gut darin ist, Informationen aufzubereiten. Natürlich ist KI auch gut darin, Code zu schreiben und zu bereinigen. Wir halten unseren Umgang mit KI für gesund, weil er ihre Stärken ausnutzt, wie unser Vorgehen zeigt.

Konkret haben wir die KI beim Brainstorming und bei der Planung nach ihrer Einschätzung gefragt. Für die Recherche haben wir selbst über Google Scholar gesucht, unsere Bücher hat uns allerdings die KI vorgeschlagen, als wir nach Material gefragt haben. Im Code haben wir KI nirgends eingesetzt, wo es für die Erkenntnisfrage relevant war, und eher für Prototypen als für fertige Systeme. Am fertigen Spiel ist KI nur bei den Shadern beteiligt, die sie für uns portiert oder selbst geschrieben hat, sowie beim Soundsystem. Beim Arbeitsjournal wurde KI genutzt, um unsere Einträge so zu strukturieren, wie es unsere Betreuungsperson verlangt hatte.

Im Begleitbericht hat die KI die Rückverweise gesetzt und ein Glossar für Begriffe erstellt, die unverständlich sein könnten. Neben dem Glossar wurden einige wenige weitere Texte mit KI formuliert; diese kann man in der KI-Deklaration (Abschnitt~\ref{sec:KIDeklaration}) nachschlagen. Zudem hat sie den gesamten Begleitbericht auf Rechtschreibung geprüft und nach logischen Fehlern gesucht, die wir über rund 100 Seiten hinweg selbst nicht im Auge behalten konnten. Die Playtest-Protokolle hat sie anonymisiert und geschlechtsneutral gemacht, und die \LaTeX-Formatierung haben wir ihr diktiert, worauf sie diese für uns umgesetzt hat.

Die KI hatte aber auch ihre Grenzen. Über die Dauer des Projekts haben wir mehrere Chatverläufe verloren, sodass der KI manchmal der Kontext fehlte und wir ihn in einem neuen Chat erneut aufbauen mussten. Zudem kann die KI weder wissen, was in unseren Köpfen vorgeht, noch den Zusammenhang des gesamten Projekts kennen, den wir kennen.

Als Fazit sehen wir die grösste Gefahr bei der Nutzung von KI darin, das Projekt an sie zu verlieren, weil der Eigenanteil zu klein und der Anteil der KI zu gross wird. Wir haben unsere Nutzung von KI zudem eingeschränkt, da wir unsere Kreativität bewahren wollten und es uns wichtig war, inhaltliche Entscheidungen selbst zu treffen. Inhaltlich war KI für uns nur ein Mittel, unsere Kreativität anzuregen.

\end{document}
